\documentclass[12pt,a4paper]{article}

\usepackage[UTF8]{ctex}
\usepackage[margin=1in]{geometry}
\usepackage{amsmath,amssymb,amsthm}
\usepackage{graphicx}
\usepackage{booktabs}
\usepackage{enumitem}
\usepackage{hyperref}
\usepackage{xcolor}
\usepackage{caption}
\usepackage{tcolorbox}
\usepackage{multirow}
\usepackage{makecell}
\usepackage{tikz}
\usetikzlibrary{arrows.meta,positioning}

% === Proposal Colors ===
\definecolor{prop5}{HTML}{E67E22}   % Orange
\definecolor{prop6}{HTML}{1ABC9C}   % Teal
\definecolor{prop7}{HTML}{E74C3C}   % Red
\definecolor{prop8}{HTML}{3498DB}   % Blue
\definecolor{prop9}{HTML}{8E44AD}   % Deep Purple
\definecolor{accent}{HTML}{2C3E50}  % Dark

% === Proposal Box ===
\newtcolorbox{proposalbox}[2][]{%
  colback=#2!5,colframe=#2!70,boxrule=0.5pt,
  fonttitle=\bfseries,title={#1},
  left=6pt,right=6pt,top=4pt,bottom=4pt
}

% === Insight Box ===
\newtcolorbox{insightbox}[1][]{%
  colback=yellow!5,colframe=yellow!60!black,boxrule=0.4pt,
  fonttitle=\bfseries,title={Core Insight},#1
}

% === Risk Box ===
\newtcolorbox{riskbox}[1][]{%
  colback=red!3,colframe=red!40,boxrule=0.3pt,
  fonttitle=\bfseries,title={Open Risks},#1
}

% === Experiment Box ===
\newtcolorbox{expbox}[1][]{%
  colback=blue!3,colframe=blue!40,boxrule=0.3pt,
  fonttitle=\bfseries,title={Key Experiments},#1
}

% === Why Box (for first-principles chains) ===
\newtcolorbox{whybox}[1][]{%
  colback=gray!5,colframe=gray!60,boxrule=0.4pt,
  fonttitle=\bfseries,title={First-Principles Chain},#1
}

\hypersetup{
  colorlinks=true,
  linkcolor=black!70,
  citecolor=blue!60!black,
  urlcolor=blue!70!black,
}

\title{\textbf{VLM Post-Training 研究方向提案}\\
\large{从第一性原理重构 Vision-Language 对齐范式}\\[6pt]
\normalsize{Internal Research Directions --- Multimodal Extension}}

\author{Jinhong Lin}
\date{March 2026}

\begin{document}

\maketitle

% ============================================================
\section*{概览}
% ============================================================

本文档提出五个独立的 VLM post-training 研究方向，
作为前文四个 text-centric 方向（统一 Reward 框架、Meta-Capability、
对比式自我诊断、多视角稳健性验证）的\textbf{多模态伴侣}。

每个方向均源自对 VLM 特有困境的第一性原理分析——
反复追问 ``why''，直到触及最深层的结构性限制。

\begin{insightbox}
Text-only post-training 的核心瓶颈是``模型不知道自己不知道什么''。
VLM post-training 的核心瓶颈更深一层：
\textbf{模型不知道自己没看到什么}。

在文本世界中，模型的全部输入在 context 中——它至少\textit{有机会}处理所有信息。
在多模态世界中，视觉信息经过编码器的有损压缩后，
原始图像中的大量细节已经在模型``看到''之前就消失了。
Post-training 在 LLM 端做的一切优化，
都受限于 visual encoding 阶段已经发生的信息损失。
\end{insightbox}

\noindent
五个方向的布局如下：
\begin{itemize}[leftmargin=2em]
  \item \textbf{方向五}（Active Re-Perception）攻击最深层限制——
    让模型在推理过程中\textit{主动}重新查看图像，
    将视觉从``被动输入''提升为``可查询记忆''。
  \item \textbf{方向六}（Modality-Decomposed Uncertainty）
    将 text-only 的 meta-cognition 框架扩展到多模态，
    揭示 VLM 特有的\textit{双重不确定性}结构。
  \item \textbf{方向七}（Cross-Modal Self-Verification）
    利用多模态的天然结构优势——
    图像本身就是 free reward signal，
    打破 reward bottleneck。
  \item \textbf{方向八}（Connector as Alignment Bottleneck）
    挑战主流范式——
    post-training 应该训练 connector 而非仅训练 LLM。
  \item \textbf{方向九}（2D Difficulty Decomposition）
    将 DPE 的 $p(1-p)$ 原则推广到二维——
    感知难度和推理难度是正交的，
    需要独立的课程设计。
\end{itemize}

\noindent
与前四个方向的关系：
\begin{itemize}[leftmargin=2em]
  \item 方向六直接扩展方向一（Self-Consistency）和方向二（Meta-Capability）到多模态。
  \item 方向七是方向四（Robustness Verification）的多模态特化——
    利用图像作为比 text reframing 更强的验证信号。
  \item 方向九是方向三（Contrastive Self-Diagnosis）的多模态推广——
    在二维空间中寻找 frontier。
  \item 方向五和方向八是\textbf{纯 VLM 原生}的问题，
    在 text-only 设定中无对应物。
\end{itemize}


\newpage
% ============================================================
\section{方向五：主动视觉回看——推理即感知}
\label{sec:prop5}
% ============================================================

\begin{proposalbox}[Active Visual Re-Perception: Turning Vision from Passive Input to Queryable Memory]{prop5}
\textbf{一句话摘要：}
当前 VLM 处理图像一次，然后在纯文本空间中推理——
就像一个人看了一眼照片，然后闭上眼睛回答所有问题。
真正的视觉推理需要\textbf{在推理过程中反复回看图像}。
\end{proposalbox}

\subsection{第一性原理推导}

\begin{whybox}
\textbf{Why} do VLMs hallucinate about visual details?\\
$\rightarrow$ 因为模型在推理过程中无法``回看''图像验证自己的描述。\\[3pt]
\textbf{Why} can't the model look back at the image?\\
$\rightarrow$ 因为图像在输入阶段被编码为固定的 visual tokens，
之后只能通过 attention 机制间接访问——
随着推理链变长，这些 visual tokens 在 attention 中的影响逐渐衰减。\\[3pt]
\textbf{Why} does visual attention decay matter?\\
$\rightarrow$ 因为长链推理（10+ 步骤）中，
最后几步的推理几乎完全依赖前面步骤的\textit{文本}总结，
而非原始视觉信息。
如果前面某步的文本总结遗漏了关键视觉细节，
后续步骤永远无法恢复这一信息。\\[3pt]
\textbf{Root cause:}
视觉信息在当前架构中是\textbf{一次性消费品}，
而非可反复查询的\textbf{持久记忆}。
\end{whybox}

\subsection{核心论点}

\begin{insightbox}
人类解视觉问题时，``看''不是发生在开头的一次性事件——
而是\textbf{贯穿整个推理过程的主动行为}。
数学家看几何题时反复查看图形，
医生读 CT 片时会在不同窗位间切换。
``看回去''不是辅助动作——它是推理的\textit{组成部分}。
\end{insightbox}

\noindent
MVoT（``在图像中思考''）迈出了重要一步——
让模型在推理过程中\textit{生成}可视化。
但它解决的是另一个问题：
生成新的视觉表征来辅助推理（类似在草稿纸上画图）。
本方向解决的是更基础的问题：
\textbf{重新审视已有的视觉输入}。

差异类比：
MVoT $\approx$ 学生在草稿纸上画辅助线；
本方向 $\approx$ 学生抬头重新看了一眼黑板上的题目图。
两者互补，但``看回去''是更基本的能力。

\subsection{方法：Active Re-Perception Training}

\paragraph{Mechanism Design.}
在 VLM 的 chain-of-thought 中引入一种新的\textbf{推理动作}：
\texttt{<look\_at region="...">}。
当模型在推理过程中发出这一 token 序列时，
系统对原始图像的指定区域进行\textbf{重新编码}，
产生新的 visual tokens 注入到当前 context 中。

\begin{tcolorbox}[colback=gray!5,colframe=gray!50,boxrule=0.3pt,left=4pt,right=4pt]
\small
\texttt{Question: 图中三角形的面积是多少？}\\[3pt]
\texttt{Step 1: 让我先整体观察这个图形...}\\
\texttt{我看到一个三角形，顶点似乎在 A, B, C 处。}\\[3pt]
\texttt{Step 2: 我需要确认底边的长度。}\\
\texttt{<look\_at region="bottom-edge">}\\
\texttt{[新的高分辨率 visual tokens 被注入]}\\
\texttt{底边标注为 8cm，端点确实是 B 和 C。}\\[3pt]
\texttt{Step 3: 现在我需要高的精确值。}\\
\texttt{<look\_at region="height-annotation">}\\
\texttt{[新的 visual tokens]}\\
\texttt{高标注为 5cm，从 A 到底边的垂足。}\\[3pt]
\texttt{Step 4: 面积 = $\frac{1}{2} \times 8 \times 5 = 20$ cm$^2$}
\end{tcolorbox}

\paragraph{Region Specification.}
\texttt{region} 可以是：
\begin{itemize}[leftmargin=2em]
  \item \textbf{Spatial coordinates}：
    normalized bounding box \texttt{[x1,y1,x2,y2]}，
    利用 Qwen2.5-VL 已具备的 grounding 能力。
  \item \textbf{Semantic reference}：
    \texttt{"the label on the x-axis"}，
    由模型的 grounding 能力解析为坐标。
  \item \textbf{Zoom level}：
    \texttt{zoom=2x}，
    对指定区域以更高分辨率重新编码。
\end{itemize}

\paragraph{Training: 三阶段。}

\textbf{Stage 1: 构造``需要回看''的数据。}

核心原则：只在模型\textit{实际因为没有回看而犯错}的地方
构造回看训练数据。
\begin{enumerate}[leftmargin=2em]
  \item 收集视觉推理问题（几何、图表分析、医学影像、地图阅读等）。
  \item 让 VLM 不带回看机制回答，收集\textbf{错误轨迹}。
  \item 分析错误：哪些错误源于\textit{视觉信息遗漏}
    （如读错数字、忽略标注、混淆相近区域）？
  \item 对这些错误，标注``如果在 Step N 回看了 Region R，
    模型可以看到 [具体信息] 并纠正错误''。
  \item 构造``回看后纠正''的 gold trajectory。
\end{enumerate}

\textbf{Stage 2: SFT on Re-Perception Trajectories.}

训练模型生成包含 \texttt{<look\_at>} 动作的推理链。
关键：同时训练模型\textbf{何时不回看}——
简单问题不需要回看，强制回看反而增加无用 token。
在 gold data 中混入``不需要回看的简单问题''的 trajectories。

\textbf{Stage 3: RL on Re-Perception Efficiency.}

SFT 后，模型可能过度使用或不足使用回看。
RL reward 设计：
\begin{itemize}[leftmargin=2em]
  \item 正确 $+$ 适度回看 $\rightarrow$ 最高 reward
  \item 正确 $+$ 无回看（不需要时）$\rightarrow$ 高 reward
  \item 正确 $+$ 过多回看 $\rightarrow$ 中等 reward（正确但浪费）
  \item 错误 $+$ 本可通过回看避免 $\rightarrow$ 强负 reward
  \item 回看后仍然错误 $\rightarrow$ 轻度负 reward（至少尝试了）
\end{itemize}

\noindent
这一 reward 结构联合优化了\textbf{正确性}和\textbf{效率}——
模型学会在需要时回看，不需要时直接推理。

\subsection{为什么不是 Tool Use？}

表面上，\texttt{<look\_at>} 与 agentic tool use 相似。
但存在本质区别：

\begin{table}[h]
\centering
\small
\begin{tabular}{lll}
\toprule
\textbf{维度} & \textbf{Agentic Tool Use} & \textbf{Active Re-Perception} \\
\midrule
延迟 & 高（API 调用） & 低（模型内部操作） \\
信息源 & 外部（搜索引擎、代码解释器） & 内部（已有的输入图像） \\
训练目标 & 学会何时调用外部工具 & 学会何时重新审视已有信息 \\
认知类比 & 查阅参考资料 & 重新看了一眼原始材料 \\
\bottomrule
\end{tabular}
\end{table}

\noindent
Active re-perception 是比 tool use 更基本的能力——
它不依赖外部系统，而是模型对自身输入的\textbf{主动管理}。
这更接近人类认知中的 \textbf{selective attention}。

\begin{expbox}
\begin{enumerate}[leftmargin=1.5em]
  \item \textbf{Hallucination reduction}：
    在 POPE、HallusionBench 上对比
    标准 VLM vs.\ 带 re-perception 的 VLM。
    假设：re-perception 显著减少 \textit{visual detail} 类幻觉
    （读错数字、忽略标注），
    对 \textit{reasoning} 类幻觉影响较小——
    证明方法精准攻击了视觉感知类错误。

  \item \textbf{Long-chain reasoning}：
    在 MathVista、AI2D 等需要多步视觉推理的 benchmark 上测试。
    假设：推理链越长，re-perception 的收益越大——
    因为 visual attention decay 是累积性的。

  \item \textbf{Efficiency analysis}：
    追踪模型平均每道题使用多少次 re-perception。
    对比 RL 前后：RL 应使回看次数下降（学会了何时不需要），
    但在困难问题上回看次数上升（学会了何时需要）。

  \item \textbf{Ablation: re-perception vs.\ higher resolution}：
    直接使用更高分辨率的初始 visual encoding 作为 baseline。
    假设：re-perception 在\textit{总 compute 相同}的情况下更优——
    因为它将 compute 分配到模型\textit{实际需要}的区域，
    而高分辨率将 compute 均匀分配到所有区域。
\end{enumerate}
\end{expbox}

\begin{riskbox}
\begin{itemize}[leftmargin=1.5em]
  \item \textbf{实现复杂度}。
    \texttt{<look\_at>} 需要在推理过程中动态注入新 token，
    这与标准 autoregressive generation 不完全兼容。
    需要修改 inference pipeline 或设计巧妙的 token 拼接方案。
  \item \textbf{Region specification 准确性}。
    模型可能在不知道该看哪里时无法准确指定 region。
    ``我不知道我不知道什么''的问题在空间域中更加尖锐。
  \item \textbf{训练数据的标注成本}。
    构造``需要回看''的 gold trajectory 需要精确标注
    ``错误发生在哪一步''和``回看哪个区域可以修复''。
    可通过强模型（GPT-4o）自动标注来缓解。
\end{itemize}
\end{riskbox}


\newpage
% ============================================================
\section{方向六：模态解耦的不确定性——我没看到 vs.\ 我没想到}
\label{sec:prop6}
% ============================================================

\begin{proposalbox}[Modality-Decomposed Uncertainty: Separating ``I Didn't See'' from ``I Didn't Think'']{prop6}
\textbf{一句话摘要：}
VLM 的错误有两个独立来源——感知错误和推理错误。
当前方法将两者混为一谈。
将 meta-cognition 解耦为\textbf{感知置信度}和\textbf{推理置信度}，
可以实现精准的错误归因和靶向训练。
\end{proposalbox}

\subsection{第一性原理推导}

\begin{whybox}
\textbf{Why} is VLM calibration worse than text-only LLM calibration?\\
$\rightarrow$ 因为 VLM 有\textbf{两个}独立的不确定性来源：
视觉感知和语言推理。\\[3pt]
\textbf{Why} does having two sources matter?\\
$\rightarrow$ 因为模型的总体置信度是两者的混合——
一个``中等置信度''可能意味着
``看得很清楚但推理不确定''或``看不清楚但推理很确定''。
这两种情况需要完全不同的应对策略。\\[3pt]
\textbf{Why} don't current methods distinguish them?\\
$\rightarrow$ 因为 DPO/RLHF 的 reward 信号是\textbf{outcome-level} 的——
``答案对了还是错了''。它不区分\textit{为什么}错了。\\[3pt]
\textbf{Root cause:}
VLM 的错误有\textbf{模态结构}（modality structure），
但训练信号是\textbf{无结构的}（unstructured）。
\end{whybox}

\subsection{核心论点}

\begin{insightbox}
考虑两个 VLM 错误案例：
\begin{enumerate}[leftmargin=2em]
  \item 模型看到图中写着``7''但读成了``1''，
    然后基于``1''做了完美的推理 $\rightarrow$ \textbf{感知错误}。
  \item 模型正确读出图中的``7''，
    但在后续数学推理中犯了逻辑错误 $\rightarrow$ \textbf{推理错误}。
\end{enumerate}
这两种错误的训练信号应该\textit{完全不同}：
案例 1 需要强化视觉感知（更细粒度的编码、更仔细的阅读），
案例 2 需要强化逻辑推理（更多验证步骤、更谨慎的推理链）。
但在当前 DPO/RLHF 范式下，两者获得\textit{相同}的负 reward——
``答案错了''。
\end{insightbox}

\subsection{方法：双维度 Meta-Cognition 训练}

\paragraph{Step 1: 构建错误归因数据集。}
对 VLM 的错误轨迹进行\textbf{模态归因}：

\begin{enumerate}[leftmargin=2em]
  \item 收集 VLM 在视觉推理任务上的错误回答。
  \item 对每个错误，构造\textbf{诊断探测}：
    \begin{itemize}
      \item \textbf{感知探测}：
        将图像中的关键信息\textit{用文本描述}替换图像输入，
        让模型重新回答。
        若模型在纯文本输入下回答正确 $\rightarrow$ 原始错误是\textbf{感知错误}。
      \item \textbf{推理探测}：
        将模型的视觉感知结果（OCR 提取、物体描述等）
        直接作为 ground truth 告诉模型，
        让模型基于这些 ground truth 重新推理。
        若模型仍然错误 $\rightarrow$ 原始错误包含\textbf{推理错误}。
    \end{itemize}
  \item 分类每个错误为：
    \begin{itemize}
      \item \textbf{P-error}（Pure Perception）：
        感知探测通过 + 推理探测通过 $\rightarrow$ 纯感知问题
      \item \textbf{R-error}（Pure Reasoning）：
        感知探测失败 + 推理探测失败 $\rightarrow$ 纯推理问题
      \item \textbf{PR-error}（Compound）：
        两个探测结果不一致 $\rightarrow$ 感知和推理的复合错误
    \end{itemize}
\end{enumerate}

\paragraph{Step 2: 构建 2$\times$2 置信度训练数据。}

类似方向二的 meta-capability training，
但要求模型输出\textbf{两维}的 self-assessment：

\begin{tcolorbox}[colback=gray!5,colframe=gray!50,boxrule=0.3pt,left=4pt,right=4pt]
\small
\texttt{Question: [附带图像的问题]}\\[3pt]
\texttt{Visual perception check:}\\
\texttt{- 我从图中读取到了: [模型的感知结果]}\\
\texttt{- 感知置信度: [high / medium / low]}\\
\texttt{- 不确定的感知元素: [如果有，列出]}\\[3pt]
\texttt{Reasoning check:}\\
\texttt{- 基于上述感知，我的推理是: [推理链]}\\
\texttt{- 推理置信度: [high / medium / low]}\\
\texttt{- 推理中最不确定的步骤: [如果有，指出]}\\[3pt]
\texttt{Overall answer: [最终答案]}\\
\texttt{Overall confidence: [基于两个维度的综合判断]}
\end{tcolorbox}

\paragraph{Step 3: RL on Decomposed Calibration.}

设计一个\textbf{结构化 reward function}：

\begin{align}
R = \underbrace{R_{\text{correct}}}_{\text{答案正确性}}
  + \lambda_P \cdot \underbrace{R_{\text{P-cal}}}_{\text{感知校准}}
  + \lambda_R \cdot \underbrace{R_{\text{R-cal}}}_{\text{推理校准}}
  + \lambda_I \cdot \underbrace{R_{\text{interact}}}_{\text{交互项}}
\end{align}

\noindent
其中：
\begin{itemize}[leftmargin=2em]
  \item $R_{\text{P-cal}}$：
    模型声称``感知置信度高''但实际是 P-error $\rightarrow$ 强负 reward。
  \item $R_{\text{R-cal}}$：
    模型声称``推理置信度高''但实际是 R-error $\rightarrow$ 强负 reward。
  \item $R_{\text{interact}}$：
    模型正确归因错误类型
    （``我的感知可能有误''当它确实是 P-error）$\rightarrow$ 正 reward。
    这是最有价值的信号——教模型\textit{知道错误的来源}。
\end{itemize}

\subsection{为什么解耦比合并更好}

\textbf{直觉论证：}
一个 pilot 的仪表板将\textit{高度}和\textit{速度}分开显示，
而非合并为一个``飞行状态''指标。
因为``飞得太低但速度正常''和``高度正常但速度太快''
需要完全不同的操作。
合并信号会导致``指标正常但即将坠机''。

\textbf{信息论论证：}
合并后的 uncertainty 是标量，
解耦后的 uncertainty 是二维向量——
后者对错误归因的信息量严格更大。
且在训练时，靶向 P-error 和 R-error 的数据可以独立更新
各自的 meta-cognitive circuit，
避免``训练推理 meta-cognition 时意外破坏感知 meta-cognition''的干扰。

\subsection{与方向一/二的关系}

方向一（Self-Consistency）的四桶框架
在 VLM 中自然扩展为\textbf{八桶}
（2 consistency $\times$ 2 correctness $\times$ 2 modality-attribution）。
但八桶可能过于细粒度，
实际上关键的区分只有三类：P-error、R-error、PR-error。

方向二（Meta-Capability）的杠杆层级在 VLM 中增加一个新的最高杠杆：
\textbf{Perception Calibration}——
知道\textit{自己是否正确感知了图像}，
是所有后续推理的前提。
这比 text-only 的 calibration 更基本，
因为文本输入的``感知''是确定的（token 就在那里），
而视觉输入的``感知''是不确定的（编码可能丢失信息）。

\begin{expbox}
\begin{enumerate}[leftmargin=1.5em]
  \item \textbf{Error attribution accuracy}：
    在 MathVista、ChartQA 上收集 VLM 错误，
    用诊断探测进行自动归因。
    报告 P-error vs.\ R-error 的比例分布。
    假设：P-error 在 fine-grained tasks（图表、文档）中占主导，
    R-error 在逻辑推理任务中占主导。
    这一分布本身就是重要的 empirical contribution。

  \item \textbf{Decomposed vs.\ unified calibration}：
    对比本方向的 2D calibration 训练
    vs.\ 方向二直接应用到 VLM（1D calibration）。
    假设：2D 版本在 ECE/Brier score 上更优，
    尤其在 mixed-type tasks（同时需要精确感知和多步推理）上
    差距最大。

  \item \textbf{Targeted training experiment}：
    识别出 P-error 和 R-error 后，
    分别构造靶向数据训练。
    对比：全量混合训练 vs.\ 按错误类型分流训练。
    假设：分流训练在\textit{总训练量更少}的情况下
    达到更高的性能。

  \item \textbf{Human preference correlation}：
    模型输出 2D confidence 后，
    用户更信任哪种模型——
    能说``我不太确定图中的数字''的模型，
    还是只能说``我不太确定''的模型？
    假设：细粒度的不确定性表达显著提升用户信任和使用效率。
\end{enumerate}
\end{expbox}

\begin{riskbox}
\begin{itemize}[leftmargin=1.5em]
  \item \textbf{诊断探测的准确性}。
    ``用文本替换图像''可能改变任务性质——
    有些问题在纯文本输入下本身变得更简单。
    需要仔细设计控制实验。
  \item \textbf{P-error 和 R-error 的纠缠}。
    实际中，很多错误是 compound 的——
    模型\textit{部分}感知正确但遗漏了关键细节，
    然后基于不完整的感知做了\textit{合理}但错误的推理。
    二分法可能过度简化。
  \item \textbf{Two-dimensional calibration 的评估指标}。
    ECE 和 Brier score 是一维的。
    需要设计新的 2D calibration metrics，
    这本身是一个非平凡的技术贡献。
\end{itemize}
\end{riskbox}


\newpage
% ============================================================
\section{方向七：跨模态自验证——图像即 Free Reward}
\label{sec:prop7}
% ============================================================

\begin{proposalbox}[Cross-Modal Self-Verification: The Image as Free Reward Signal]{prop7}
\textbf{一句话摘要：}
在多模态场景中，你拥有 text-only 所没有的天然验证信号——
图像本身。
通过检查模型的文本输出与原始图像之间的\textbf{跨模态一致性}，
可以构造不依赖任何人类标注的 reward signal，
打破 VLM alignment 的 reward bottleneck。
\end{proposalbox}

\subsection{第一性原理推导}

\begin{whybox}
\textbf{Why} is VLM alignment bottlenecked by reward quality?\\
$\rightarrow$ 因为评估``文本描述是否忠实于图像''需要多模态理解——
而 reward model 本身就是一个 VLM，继承了 VLM 的所有缺陷。\\[3pt]
\textbf{Why} can't we build a perfect multimodal reward model?\\
$\rightarrow$ 因为它需要大量人类标注的``好/坏''偏好对，
而多模态偏好标注的成本是 text-only 的 5--10 倍
（标注者需要仔细对比文本与图像）。\\[3pt]
\textbf{Why} do we need an external reward model at all?\\
$\rightarrow$ 问题的反面：\textbf{图像本身就是 ground truth}。
当模型声称``图中有一只红色的猫''，
图像\textit{就在那里}——
我们不需要人类来判断这句话对不对，
我们需要的是一种\textit{自动检查}文本-图像一致性的机制。\\[3pt]
\textbf{Root cause:}
多模态场景中存在 text-only 不具备的天然验证信号——
\textbf{原始图像}。
当前方法完全没有利用这一信号。
\end{whybox}

\subsection{核心论点}

\begin{insightbox}
在 text-only 设定中，验证模型输出需要外部裁判——
因为没有独立的``真相源''。
在多模态设定中，图像\textit{就是}真相源。
模型的文本输出要么忠实于图像，要么不忠实——
这是客观可检查的。

关键洞察：我们不需要\textbf{一个新的 reward model}来检查一致性。
我们可以让\textbf{同一个 VLM}通过结构化的交叉验证来自我检查。
这不是 self-consistency（对同一问题采样多次），
而是 \textbf{cross-modal consistency}——
利用多模态的\textit{结构}来创造新的验证路径。
\end{insightbox}

\subsection{方法：三层跨模态验证}

\paragraph{Layer 1: Description $\leftrightarrow$ Image 一致性。}

最直接的验证：模型描述了图像，
然后基于\textit{仅描述}回答关于图像的问题，
再基于\textit{原始图像}回答同一问题。
两组答案是否一致？

\begin{tcolorbox}[colback=gray!5,colframe=gray!50,boxrule=0.3pt,left=4pt,right=4pt]
\small
\textbf{Step 1:} 模型描述图像 $\rightarrow$ 得到 Description $D$\\[2pt]
\textbf{Step 2:} 从 $D$ 中提取 $N$ 个可验证的事实声明
（``图中有 3 个人''、``背景是蓝色的''、``最左边的人穿红衣服''）\\[2pt]
\textbf{Step 3:} 对每个声明，基于\textit{原始图像}
向模型提问：``图中有几个人？''``背景是什么颜色？''\\[2pt]
\textbf{Step 4:} 检查 Step 3 的回答与 Step 2 的声明是否一致。
不一致 $\rightarrow$ Description 中存在幻觉。
\end{tcolorbox}

\paragraph{Layer 2: Counterfactual 一致性。}

更强的验证：构造反事实问题。

\begin{tcolorbox}[colback=gray!5,colframe=gray!50,boxrule=0.3pt,left=4pt,right=4pt]
\small
如果模型声称``图中有一只猫在沙发上''：\\[2pt]
\textbf{Counterfactual Q1:}
``如果图中没有猫，沙发上是什么？''\\
$\rightarrow$ 模型应该说``但图中确实有猫''或给出与原始一致的描述。\\[2pt]
\textbf{Counterfactual Q2:}
``图中的猫是在桌子上还是地板上？''\\
$\rightarrow$ 如果模型声称的``沙发上''是正确的，
它应该拒绝这两个选项。
如果它接受了某个选项 $\rightarrow$ 原始声明可能不可靠。
\end{tcolorbox}

\paragraph{Layer 3: Cross-Task 一致性。}

最深层的验证：同一图像的不同任务应产生一致的视觉理解。

\begin{tcolorbox}[colback=gray!5,colframe=gray!50,boxrule=0.3pt,left=4pt,right=4pt]
\small
\textbf{Task A:} ``描述这张图片''
$\rightarrow$ ``一个穿蓝色衬衫的男人在公园里跑步''\\[2pt]
\textbf{Task B:} ``图中人物在做什么运动？'' $\rightarrow$ ``瑜伽''\\[2pt]
$\rightarrow$ 不一致！Task A 说跑步，Task B 说瑜伽。\\
$\rightarrow$ 至少一个回答包含幻觉。\\
$\rightarrow$ 可以进一步通过更多任务来三角定位哪个是错的。
\end{tcolorbox}

\paragraph{Reward Signal Construction.}

将三层验证的结果聚合为 reward signal：

\begin{align}
R_{\text{cross-modal}} = \alpha \cdot \text{L1-score} + \beta \cdot \text{L2-score} + \gamma \cdot \text{L3-score}
\end{align}

\noindent
其中每个 score $\in [-1, 1]$，
一致 $\rightarrow$ 正，不一致 $\rightarrow$ 负。
这个 reward 可以直接用于 DPO（选择跨模态一致性更高的 response）
或 RL（作为在线 reward）。

\subsection{为什么优于现有方法}

\begin{table}[h]
\centering
\small
\begin{tabular}{p{3.5cm}p{3cm}p{3cm}p{3cm}}
\toprule
\textbf{方法} & \textbf{Reward 来源} & \textbf{成本} & \textbf{质量天花板} \\
\midrule
LLaVA-RLHF
  & 人类偏好
  & 极高
  & 人类标注者的视觉理解 \\
RLHF-V
  & 细粒度人类反馈
  & 非常高
  & 单步标注质量 \\
RLAIF-V
  & AI 反馈
  & 中等
  & Teacher VLM 的能力 \\
\textbf{本方向}
  & 图像自身
  & \textbf{零人工成本}
  & 模型跨模态验证能力 \\
\bottomrule
\end{tabular}
\end{table}

\noindent
唯一的``成本''是推理阶段的额外 compute（生成验证问题和答案）。
但这比人类标注便宜几个数量级。

\subsection{与方向四的关系}

方向四（Robustness Verification）在 text-only 中通过 semantic reframing
来 stress-test 推理。
本方向在多模态中有一个更强的验证信号：
\textbf{图像提供了 reframing 无法提供的客观 ground truth}。

Text reframing 只能检查``一致性''，无法检查``正确性''——
模型可能在所有 framing 下一致地错。
但 cross-modal verification 可以检查``忠实性''——
因为图像\textit{就在那里}，
任何与图像矛盾的描述都是客观可检测的幻觉。

\begin{expbox}
\begin{enumerate}[leftmargin=1.5em]
  \item \textbf{Cross-modal reward vs.\ human reward}：
    在 LLaVA-Bench、MMHalBench 上对比
    本方向的 cross-modal reward 训练
    vs.\ 人类偏好 DPO（LLaVA-RLHF style）。
    假设：两者在 hallucination reduction 上效果相当，
    但本方向\textbf{零人工标注成本}。

  \item \textbf{Scalability test}：
    固定训练 compute，变化 reward 来源：
    1K 人类偏好对 vs.\ 100K 自动生成的 cross-modal 验证信号。
    假设：数量可以弥补质量——
    大量弱信号的累积效果匹配少量强信号。

  \item \textbf{Layer ablation}：
    L1-only vs.\ L1+L2 vs.\ L1+L2+L3。
    假设：每增加一层提供边际收益，
    但 L1 已捕获大部分幻觉信号。
    L2（counterfactual）对\textit{subtle} 幻觉
    （模型描述大致正确但细节有误）最有价值。

  \item \textbf{与可验证 reward 的校准}：
    在 ChartQA、DocVQA 等有 ground truth 的任务上，
    测量 cross-modal reward 与 ground-truth reward 的相关性。
    假设：高度正相关（$\rho > 0.7$）——
    证明 cross-modal 信号是可靠的 proxy。
\end{enumerate}
\end{expbox}

\begin{riskbox}
\begin{itemize}[leftmargin=1.5em]
  \item \textbf{自我验证的 bootstrap 问题}。
    用一个有幻觉的 VLM 来验证自己的幻觉——
    如果模型系统性地``看不到''某个物体，
    验证阶段也会``看不到''，导致假阳性。
    缓解策略：使用不同的 prompt 或 temperature 进行验证，
    增加随机性以打破系统性偏差。
  \item \textbf{验证问题的生成质量}。
    从描述中提取可验证声明并构造反事实问题
    本身就需要较强的语言能力。
    可通过模板化来简化。
  \item \textbf{计算开销}。
    三层验证意味着每个训练样本需要
    多轮额外推理（生成验证问题、回答验证问题、比较结果）。
    需要在 reward 质量和 compute 之间权衡。
\end{itemize}
\end{riskbox}


\newpage
% ============================================================
\section{方向八：Connector 是 Alignment 的真正瓶颈}
\label{sec:prop8}
% ============================================================

\begin{proposalbox}[The Connector as the Alignment Bottleneck: Post-Training the Wrong Component]{prop8}
\textbf{一句话摘要：}
当前 VLM post-training 的主流范式是``冻结视觉编码器和 connector，
只训练 LLM''。
但如果视觉信息在 connector 层就已经被有损压缩，
LLM 端的 alignment 永远无法恢复丢失的信息——
\textbf{你无法对齐你看不到的东西}。
\end{proposalbox}

\subsection{第一性原理推导}

\begin{whybox}
\textbf{Why} does VLM hallucination persist even after strong DPO/RLHF?\\
$\rightarrow$ 因为模型``看到''的视觉信息就是不完整的。
DPO 可以教模型``不要编造''，
但无法教模型``看到你没看到的东西''。\\[3pt]
\textbf{Why} doesn't the model see complete visual information?\\
$\rightarrow$ 因为 visual connector（MLP projection 或 Q-former）
将 vision encoder 输出的高维特征压缩到 LLM 的 embedding space 中，
这是一个\textit{有损}映射。\\[3pt]
\textbf{Why} is the connector lossy?\\
$\rightarrow$ 因为 connector 在 pre-training 阶段被训练——
那时的训练信号是 image-caption matching，
而非 fine-grained visual reasoning。
Connector 学会了保留对 caption 有用的信息
（粗粒度语义），
而丢弃了对 reasoning 有用但对 caption 无关的信息
（空间关系、精确数值、微妙细节）。\\[3pt]
\textbf{Root cause:}
Connector 的\textbf{信息保留策略}在 pre-training 阶段已经固化，
而 post-training 在它的下游操作——
\textbf{优化了错误的组件}。
\end{whybox}

\subsection{核心论点}

\begin{insightbox}
VLM 的 post-training pipeline 可以类比为一条\textbf{信息管道}：

$$\text{Image} \xrightarrow{\text{Vision Encoder}} \text{Features}
\xrightarrow{\text{Connector}} \text{LLM tokens}
\xrightarrow{\text{LLM}} \text{Response}$$

当前 post-training（SFT/DPO/RL）只优化管道的\textit{最后一段}（LLM）。
但信息在\textit{中间段}（Connector）就已经被压缩。
这就像你试图通过优化扬声器来改善音质——
但问题出在录音阶段的麦克风。
\end{insightbox}

\noindent
\textbf{经验证据：}
\begin{itemize}[leftmargin=2em]
  \item Cambrian-1 发现 visual representation 是 VLM 性能的关键瓶颈——
    不同 encoder 的选择比 post-training 方法的选择影响更大。
  \item InternVL2.5 发现对 vision encoder 的投资可以替代对数据的投资。
  \item 但这些工作的结论都是``换一个更好的 encoder''——
    \textbf{没有}任何工作研究过``让 alignment 信号回传到 connector''。
\end{itemize}

\subsection{方法：Alignment-Aware Connector Training}

\paragraph{核心思路。}
让 DPO/RL 的 alignment 梯度\textbf{穿过 LLM 回传到 connector}，
使 connector 学会``alignment 需要什么信息''。

\paragraph{Phase 1: 诊断——connector 丢了什么？}

\begin{enumerate}[leftmargin=2em]
  \item 收集 VLM 在细粒度视觉任务上的系统性错误。
  \item 用 probing 技术检查：
    这些错误是否源于 connector 输出中\textit{缺失}的信息？
    \begin{itemize}
      \item 训练一个 linear probe，
        尝试从 connector 输出的 visual tokens 中
        恢复 fine-grained 信息（空间坐标、精确数值、物体属性）。
      \item 如果 probe 无法恢复 $\rightarrow$ 信息在 connector 层就丢了。
      \item 如果 probe 可以恢复但 LLM 没有利用
        $\rightarrow$ 信息存在但 LLM 不知道如何提取。
    \end{itemize}
  \item 这一诊断区分了两种不同的瓶颈：
    \textbf{information loss}（connector 问题）
    vs.\ \textbf{information utilization}（LLM 问题）。
\end{enumerate}

\paragraph{Phase 2: Connector-Included Alignment Training.}

对于诊断为 information loss 的错误类型：

\begin{enumerate}[leftmargin=2em]
  \item \textbf{Unfreeze connector during DPO/RL。}
    标准做法是冻结 connector，只训练 LLM。
    本方向的核心改变：
    让 alignment loss 的梯度同时更新 connector 参数。
  \item \textbf{Selective unfreezing。}
    完全解冻 connector 可能导致 catastrophic forgetting
    （connector 丢失 pre-training 学到的基本语义对齐）。
    解决方案：
    \begin{itemize}
      \item \textbf{LoRA on connector}：
        冻结 connector 的原始权重，
        添加 low-rank adaptation 专门接受 alignment 梯度。
      \item \textbf{Residual connection}：
        在 connector 的原始输出之上添加一个可训练的 residual path，
        让 alignment 信号补充（而非替代）原始映射。
    \end{itemize}
  \item \textbf{Alignment-targeted reward。}
    在标准 DPO reward（哪个 response 更好？）之上，
    增加\textbf{visual grounding reward}——
    response 中引用的视觉细节是否确实存在于图像中？
    这一 reward 天然需要 connector 提供更完整的视觉信息，
    从而驱动 connector 的更新方向。
\end{enumerate}

\paragraph{Phase 3: 验证信息恢复。}

训练后，重复 Phase 1 的 probing 实验。
关键假设：
\begin{itemize}[leftmargin=2em]
  \item 经过 alignment-aware training 的 connector
    输出中包含更多 fine-grained 信息。
  \item 这些新保留的信息直接对应于 alignment 需要但之前被丢弃的维度。
\end{itemize}

\subsection{与``换更好的 encoder''的区别}

InternVL 和 Cambrian-1 的路线是``用更好的 encoder 替换旧的''——
这需要重新 pre-train。
本方向的路线是``让现有 connector 在 post-training 阶段
学会保留 alignment 需要的信息''——
这在任何已有 VLM 上都可以低成本实施。

类比：
``换更好的 encoder'' $\approx$ 换一台更好的相机。
``alignment-aware connector training'' $\approx$
教摄影师针对你需要的主题选择更好的取景角度——
不换相机，但拍出更好的照片。

\begin{expbox}
\begin{enumerate}[leftmargin=1.5em]
  \item \textbf{Probing experiment}（Phase 1 诊断）：
    在 LLaVA-1.5、InternVL2 上进行 connector output probing。
    测量 fine-grained 信息（spatial, numerical, attribute）
    在 connector 前后的保留率。
    \textbf{这一实验本身就是有价值的 empirical contribution}——
    定量回答``connector 到底丢了什么？''。

  \item \textbf{Connector-included DPO vs.\ standard DPO}：
    相同的 preference data，
    唯一区别是 connector 是否参与训练。
    假设：connector-included 在 fine-grained 任务
    （ChartQA、DocVQA、TextVQA）上显著更优，
    在 coarse-grained 任务（VQAv2、GQA）上差距较小——
    因为粗粒度信息 connector 本来就保留了。

  \item \textbf{LoRA vs.\ full unfreeze vs.\ residual}：
    三种 connector 参与方式的对比。
    假设：residual connection 是最佳平衡——
    保留原始对齐能力 + 补充 fine-grained 信息。
    Full unfreeze 有 catastrophic forgetting 风险。

  \item \textbf{Information recovery visualization}：
    训练前后的 connector output 的 t-SNE 可视化。
    假设：训练后的 connector 输出在
    fine-grained attribute 维度上有更好的可分性。
\end{enumerate}
\end{expbox}

\begin{riskbox}
\begin{itemize}[leftmargin=1.5em]
  \item \textbf{Connector 的 capacity 限制}。
    MLP connector 的参数量远小于 LLM。
    即使解冻，它可能没有足够的容量
    来同时保留语义对齐和 fine-grained 信息。
    可能需要增大 connector 的参数量，
    这引入了架构修改。
  \item \textbf{Vision encoder 的天花板}。
    如果 vision encoder 本身就没有编码某些信息
    （如 14$\times$14 分辨率下的微小文字），
    connector 再怎么优化也恢复不了。
    本方向的适用范围是
    ``encoder 看到了但 connector 丢了''的情况。
  \item \textbf{训练稳定性}。
    同时训练 connector 和 LLM 可能导致训练不稳定——
    两个组件的 learning rate 和 gradient scale 差异巨大。
    需要仔细的 learning rate scheduling。
\end{itemize}
\end{riskbox}


\newpage
% ============================================================
\section{方向九：二维难度分解——感知课程与推理课程的正交设计}
\label{sec:prop9}
% ============================================================

\begin{proposalbox}[2D Difficulty Decomposition: Orthogonal Curricula for Perception and Reasoning]{prop9}
\textbf{一句话摘要：}
VLM 的失败有两个正交维度：感知难度和推理难度。
当前方法将所有失败混在一起训练。
分解为二维课程，可以实现``用更少的数据达到更好的效果''——
因为每份训练数据都精准靶向\textit{实际}的瓶颈维度。
\end{proposalbox}

\subsection{第一性原理推导}

\begin{whybox}
\textbf{Why} does increasing training data for VLMs yield diminishing returns?\\
$\rightarrow$ 因为很大一部分训练数据没有命中模型的\textit{实际}弱点。
一道``模型因为看不清图而答错''的题，
如果被用来训练推理能力，就是浪费。\\[3pt]
\textbf{Why} does this mismatch happen?\\
$\rightarrow$ 因为 VLM 的 failure mode 有\textit{两个}独立来源
（感知 vs.\ 推理），
但训练数据的难度只有\textit{一个}维度
（``模型答对的概率''）。\\[3pt]
\textbf{Why} is one-dimensional difficulty insufficient?\\
$\rightarrow$ 因为``$p = 0.5$''（DPE 的 sweet spot）可以来自
完全不同的原因组合：
(a) 感知容易 + 推理困难，
(b) 感知困难 + 推理容易，
(c) 两者都中等。
这三种情况需要完全不同的训练策略。\\[3pt]
\textbf{Root cause:}
DPE 的 $p(1-p)$ 原则在 text-only 中是一维最优的，
但在多模态中退化为\textbf{低效}——
因为它将两个正交的难度维度坍缩为一个标量。
\end{whybox}

\subsection{核心论点}

\begin{insightbox}
将 DPE 的一维 difficulty landscape 推广到二维：

\begin{center}
\begin{tabular}{c|c|c|c}
 & \textbf{感知 Easy} & \textbf{感知 Medium} & \textbf{感知 Hard} \\
\hline
\textbf{推理 Easy} & 已掌握 & 感知瓶颈 & 感知瓶颈 \\
\hline
\textbf{推理 Medium} & 推理瓶颈 & \textbf{双重前沿} & 混合训练 \\
\hline
\textbf{推理 Hard} & 推理瓶颈 & 混合训练 & 超出能力 \\
\end{tabular}
\end{center}

\noindent
一维 DPE 将整个矩阵压缩成一个$p$值。
二维分解让我们为每个区域设计\textit{不同的训练策略}——
这是信息论上严格更优的课程设计。
\end{insightbox}

\subsection{方法：正交探测与分流训练}

\paragraph{Step 1: 构建正交探测框架。}

对每个视觉推理问题，构造两类探测来分离感知难度和推理难度：

\textbf{推理难度探测}（控制感知）：
\begin{itemize}[leftmargin=2em]
  \item 将图像中的关键信息替换为完美的文本描述
    （即``假装 OCR 完美、物体识别完美''），
    让模型基于\textit{完美感知}回答问题。
  \item 模型在此设定下的成功率 = 推理能力的纯度量。
\end{itemize}

\textbf{感知难度探测}（控制推理）：
\begin{itemize}[leftmargin=2em]
  \item 设计仅需\textit{简单推理}但需要\textit{精确感知}的子问题：
    ``图中 X 轴的最大值是多少？''
    ``第三行第二列的数字是什么？''
  \item 模型在此设定下的成功率 = 感知能力的纯度量。
\end{itemize}

\noindent
两个探测的交叉给出每个样本在 2D 空间中的位置。

\paragraph{Step 2: 定义四个训练区域。}

\begin{table}[h]
\centering
\small
\begin{tabular}{p{3.5cm}p{2.5cm}p{5cm}l}
\toprule
\textbf{区域} & \textbf{坐标} & \textbf{训练策略} & \textbf{权重} \\
\midrule
Mastered
  & P-easy, R-easy
  & 少量保持能力的数据
  & 5\% \\
Perception Frontier
  & P-medium, R-easy
  & 高分辨率视觉数据 + 精确 OCR/计数训练
  & 30\% \\
Reasoning Frontier
  & P-easy, R-medium
  & 长链推理数据 + step-by-step verification
  & 30\% \\
Dual Frontier
  & P-medium, R-medium
  & 综合训练 + 方向六的解耦 self-assessment
  & 25\% \\
Beyond Reach
  & P-hard, R-hard
  & 训练诚实拒绝（``这超出我的能力''）
  & 10\% \\
\bottomrule
\end{tabular}
\caption{2D difficulty decomposition 的五区域训练策略。
每个区域的训练数据\textit{格式}不同——
Perception Frontier 强调视觉细节验证步骤，
Reasoning Frontier 强调逻辑推理深度。}
\label{tab:2d-regions}
\end{table}

\paragraph{Step 3: 区域特异的数据格式。}

\textbf{Perception Frontier 数据格式}（重点是教模型仔细看）：
\begin{tcolorbox}[colback=gray!5,colframe=gray!50,boxrule=0.3pt,left=4pt,right=4pt]
\small
\texttt{Question: [需要精确视觉感知的问题]}\\[2pt]
\texttt{Visual inventory: 让我仔细列出图中的关键视觉元素...}\\
\texttt{- 我看到 [元素1]: [精确描述，包括位置、大小、颜色]}\\
\texttt{- 我看到 [元素2]: [精确描述]}\\
\texttt{- 需要双重检查的元素: [可能看错的部分]}\\[2pt]
\texttt{Based on verified visual information: [简短推理 + 答案]}
\end{tcolorbox}

\textbf{Reasoning Frontier 数据格式}（重点是教模型深入推理）：
\begin{tcolorbox}[colback=gray!5,colframe=gray!50,boxrule=0.3pt,left=4pt,right=4pt]
\small
\texttt{Question: [视觉信息清晰但推理复杂的问题]}\\[2pt]
\texttt{Visual facts (straightforward):}\\
\texttt{- [快速准确的视觉提取]}\\[2pt]
\texttt{Reasoning chain:}\\
\texttt{Step 1: [推理步骤] ... 验证: [验证逻辑]}\\
\texttt{Step 2: [推理步骤] ... 验证: [验证逻辑]}\\
\texttt{...}\\
\texttt{Answer: [最终答案]}
\end{tcolorbox}

\paragraph{Step 4: 迭代探测与更新。}

训练后重新运行正交探测，
2D landscape 会发生变化：
\begin{itemize}[leftmargin=2em]
  \item Perception Frontier 的问题可能移入 Mastered
    （感知能力提升了）。
  \item 新的问题进入 Perception Frontier
    （更难的感知挑战成为新的前沿）。
  \item 同理 Reasoning Frontier。
\end{itemize}

\noindent
这实现了\textbf{二维自适应课程}——
不再是 DPE 的一维 frontier，
而是一个在 perception-reasoning 平面上移动的 2D frontier。

\subsection{为什么二维严格优于一维}

\textbf{信息论论证：}
一维 $p$ 值是 $(\text{perception difficulty}, \text{reasoning difficulty})$
二维信号的\textit{边际化}。
边际化总是丢失信息。
在二维空间中选择训练数据，
等价于在更丰富的特征空间中做课程设计——
理论上不可能比一维更差。

\textbf{实践论证：}
考虑一个模型，它的感知能力很强但推理能力一般。
一维 DPE 会选择 $p \approx 0.5$ 的问题——
其中很多是``感知 Hard + 推理 Easy''的题。
这些题对这个模型来说，
``$p \approx 0.5$''纯粹因为感知困难——
但模型的\textit{真正}瓶颈是推理。
二维分解可以避免这种错误分配。

\subsection{与方向三/六的关系}

方向三（Contrastive Self-Diagnosis）的``失败 $\rightarrow$ 分析 $\rightarrow$ 纠正''框架
自然适配本方向：
\begin{itemize}[leftmargin=2em]
  \item Perception Frontier 的失败轨迹中，
    ``Error analysis'' 聚焦于``我在哪里看错了''。
  \item Reasoning Frontier 的失败轨迹中，
    ``Error analysis'' 聚焦于``我在哪一步推错了''。
\end{itemize}

\noindent
方向六（Decomposed Uncertainty）的 P-error / R-error 分类
直接为本方向的 2D 探测提供信号。
两者的区别：方向六关注\textit{训练后}的 runtime meta-cognition，
本方向关注\textit{训练时}的课程设计。
它们是同一洞察（``VLM 的难度是二维的''）的两个面向。

\begin{expbox}
\begin{enumerate}[leftmargin=1.5em]
  \item \textbf{2D vs.\ 1D curriculum}：
    固定总训练 compute，对比
    2D difficulty decomposition vs.\ 标准 DPE（1D）
    vs.\ random sampling。
    评估在 mixed benchmarks 上的综合表现。
    假设：2D $>$ 1D $>$ random，
    且差距在\textit{不平衡}模型上最大
    （感知 $\gg$ 推理 或 推理 $\gg$ 感知）。

  \item \textbf{Perception-Reasoning distribution map}：
    对 MathVista、ChartQA、MMMU 等 benchmark 中的问题
    绘制 2D 难度地图。
    \textbf{这一 empirical contribution 本身就有价值}——
    定量回答``这些 benchmark 到底在测感知还是推理？''。

  \item \textbf{Targeted training efficiency}：
    识别一个模型的主要瓶颈维度（如感知），
    仅在该维度上投入训练资源。
    对比``仅 perception training''的效率
    vs.\ 全量混合训练。
    假设：靶向训练以 30\% 的数据
    达到全量训练 80\% 的效果。

  \item \textbf{跨模型泛化}：
    在 LLaVA-1.5（connector = MLP）
    和 InstructBLIP（connector = Q-former）上
    分别进行 2D 探测。
    假设：不同架构的 2D profile 不同——
    MLP connector 可能在 perception 维度更弱，
    Q-former 在某些细粒度任务上更强。
    这揭示了架构选择对 difficulty landscape 的影响。
\end{enumerate}
\end{expbox}

\begin{riskbox}
\begin{itemize}[leftmargin=1.5em]
  \item \textbf{正交探测的准确性}。
    ``用文本替换图像''不是完美的感知控制——
    有些视觉信息难以用文本精确描述。
    需要设计更精细的探测方法。
  \item \textbf{维度间的交互}。
    现实中感知和推理可能不是完全正交的——
    例如，某些高级推理需要对视觉信息的\textit{精确}编码
    （不仅仅是``看到''了，还需要``精确记住''）。
    2D 分解可能遗漏这种交互效应。
  \item \textbf{数据格式设计的敏感性}。
    Perception Frontier 和 Reasoning Frontier 的
    训练数据格式差异可能引入 confounding variable。
    需要控制实验来隔离
    ``2D 选择''和``格式差异''各自的贡献。
\end{itemize}
\end{riskbox}


\newpage
% ============================================================
\section*{五个方向的统一图景}
% ============================================================

五个方向围绕一个核心问题展开：
\textbf{VLM post-training 的信息流从图像到文本，
每个环节都有独特的失败模式和优化机会}。

\begin{center}
\begin{tikzpicture}[
  node distance=0.7cm,
  box/.style={rectangle, draw, rounded corners=3pt,
    minimum width=11.5cm, minimum height=1.2cm,
    text width=11cm, align=center, font=\small},
  arrow/.style={-{Stealth[length=3mm]}, thick, gray!70}
]
\node[box, fill=prop8!10, draw=prop8!70] (p8)
  {\textbf{方向八：Connector 是 Alignment 瓶颈}\\
   攻击信息流的最上游——让 connector 保留 alignment 所需的视觉信息};
\node[box, fill=prop5!10, draw=prop5!70, below=of p8] (p5)
  {\textbf{方向五：主动视觉回看}\\
   在推理过程中重新查询图像——将视觉从被动输入变为可查询记忆};
\node[box, fill=prop6!10, draw=prop6!70, below=of p5] (p6)
  {\textbf{方向六：模态解耦不确定性}\\
   区分感知错误和推理错误——精准归因才能精准训练};
\node[box, fill=prop7!10, draw=prop7!70, below=of p6] (p7)
  {\textbf{方向七：跨模态自验证}\\
   利用图像本身作为 free reward——打破多模态 reward bottleneck};
\node[box, fill=prop9!10, draw=prop9!70, below=of p7] (p9)
  {\textbf{方向九：二维难度分解}\\
   感知课程和推理课程正交设计——用更少数据达到更好效果};

\draw[arrow] (p8) -- (p5);
\draw[arrow] (p5) -- (p6);
\draw[arrow] (p6) -- (p7);
\draw[arrow] (p6) -- (p9);
\end{tikzpicture}
\end{center}

\noindent
逻辑递进关系：
\begin{itemize}[leftmargin=2em]
  \item 方向八 $\rightarrow$ 方向五：
    先确保 connector 传递了足够的信息，
    再赋予模型主动回看的能力。
  \item 方向五 $\rightarrow$ 方向六：
    有了回看能力后，
    模型需要知道\textit{何时}回看（感知不确定时）
    vs.\ \textit{何时}深入推理（推理不确定时）。
  \item 方向六 $\rightarrow$ 方向七/九：
    有了解耦的不确定性后，
    可以用图像作为 reward（方向七）或设计正交课程（方向九）。
\end{itemize}

\subsection*{与前四个方向（Text-Only）的映射}

\begin{table}[h]
\centering
\small
\begin{tabular}{p{5cm}p{5.5cm}p{2cm}}
\toprule
\textbf{Text-Only 方向} & \textbf{VLM 对应} & \textbf{关系} \\
\midrule
方向一：统一 Self-Consistency Reward
  & 方向七：Cross-Modal Self-Verification
  & 深化 \\
方向二：Meta-Capability 假说
  & 方向六：Modality-Decomposed Uncertainty
  & 扩展 \\
方向三：对比式自我诊断
  & 方向九：2D Difficulty Decomposition
  & 推广 \\
方向四：多视角稳健性验证
  & 方向七：Cross-Modal Self-Verification
  & 特化 \\
（无对应）
  & 方向五：Active Re-Perception
  & 原生 \\
（无对应）
  & 方向八：Connector Bottleneck
  & 原生 \\
\bottomrule
\end{tabular}
\end{table}

\subsection*{优先排序}

\begin{itemize}[leftmargin=2em]
  \item \textbf{最高 risk/reward}：方向八（Connector Bottleneck）——
    若验证``connector 是真正瓶颈''，将改变整个 VLM post-training 范式。
  \item \textbf{最可行的近期实验}：方向九（2D Difficulty Decomposition）——
    仅需标准 VLM + 精心设计的探测实验，
    ``2D difficulty map'' 本身就是高价值的 empirical contribution。
  \item \textbf{最具原创性}：方向五（Active Re-Perception）——
    提出了一种全新的 VLM 推理范式，
    如果成功则是 paradigm shift。
  \item \textbf{最实用的即时价值}：方向七（Cross-Modal Self-Verification）——
    零标注成本的 reward signal，
    可以直接应用到任何 VLM alignment pipeline。
  \item \textbf{理论贡献最大}：方向六（Modality-Decomposed Uncertainty）——
    将 meta-cognition 理论从 1D 推广到结构化的 2D，
    为 VLM 的 calibration 研究奠定理论基础。
\end{itemize}

\end{document}
